\documentclass[conference]{IEEEtran}
\usepackage{amsmath,amssymb,graphicx,url,array}
\raggedbottom
\makeatletter\renewenvironment{table}[1][]{\par\noindent\begin{minipage}{\linewidth}\def\@captype{table}\footnotesize}{\end{minipage}\par}\makeatother
\begin{document}
\title{Speculative Decoding for Lossless LLM Inference Acceleration}
\author{\IEEEauthorblockN{Anonymous}}
\maketitle
\begin{abstract}
Speculative sampling accelerates LLM decoding by introducing a draft model to generate candidate tokens scored by the target model via modified rejection sampling, preserving target distribution while achieving speedups without model modifications. \cite{chen2023accelerating} The paper proposes speculative decoding as a method to accelerate inference from large autoregressive models without altering output distributions. \cite{leviathan2022fast}
\end{abstract}
\begin{IEEEkeywords}speculative decoding, draft model, lossless acceleration\end{IEEEkeywords}
\section{Introduction}
Chen et al. propose speculative sampling to accelerate LLM decoding via draft model candidate generation and target model scoring. \cite{chen2023accelerating}

Leviathan, Kalman, and Matias propose speculative decoding to accelerate inference from large autoregressive Transformers. \cite{leviathan2022fast}

\section{Reviewed sources}
\subsection{Accelerating Large Language Model Decoding with Speculative Sampling}
Speculative sampling introduces a draft model to generate short candidate drafts scored by a target model using modified rejection sampling, enabling multiple token acceptance per target call and achieving 2-2.5x decoding speedup on distributed Chinchilla 70B without model modifications. \cite{chen2023accelerating}

\subsection{Fast Inference from Transformers via Speculative Decoding}
Leviathan et al. introduce speculative decoding to accelerate inference of large autoregressive Transformers. The method uses an efficient approximation model to generate draft tokens verified in parallel against the target model, enabling exact decoding without retraining or architectural changes. \cite{leviathan2022fast}

Sun et al. describe TriForce, a hierarchical speculative decoding system that uses retrieval-based drafting with a lightweight model and StreamingLLM cache to accelerate long-context generation. Its two-phase approach leverages attention sparsity and contextual locality for high acceptance rates. \cite{sun2024triforce}

\section{Source comparison}
\begin{table}[H]\centering
\caption{Qualitative comparison of the processed sources.}
\begin{tabular}{|>{\raggedright\arraybackslash}p{0.26\linewidth}|>{\raggedright\arraybackslash}p{0.235\linewidth}|>{\raggedright\arraybackslash}p{0.235\linewidth}|}\hline
Source & Approach & Evaluation \\ \hline
Chen 2023 \cite{chen2023accelerating} & Speculative sampling with modified rejection sampling & XSum and HumanEval \\ \hline
Leviathan 2022 \cite{leviathan2022fast} & speculative decoding & T5-XXL translation and summarization \\ \hline
\end{tabular}\end{table}
\section{Evidence limitations}
The source authors state: "Note that as \ensuremath{\mathit{K}}
increases, the overall speedup plateaus or even regresses, with XSum being optimal at \ensuremath{\mathit{K}} = 3." \cite{chen2023accelerating}

The source authors state: "One limitation of speculative execution in general, and of
speculative decoding in particular, is that latency is improved through increased concurrency at the cost of an increased number of arithmetic operations." \cite{leviathan2022fast}

\section{Conclusion}
The algorithm relies on the observation that latency of parallel scoring of short continuations generated by a faster draft model is comparable to sampling a single token from the larger target model. \cite{chen2023accelerating}

The method provides speedups when memory bandwidth is the bottleneck with no output distribution changes. \cite{leviathan2022fast}

\bibliographystyle{IEEEtran}
\bibliography{refs}\end{document}
